\section{Step 31 Add for the current multiplier bit}\label{rsa-s31}

If the whole private multiplier is zero, go to publication. Otherwise test its low bit. For a one bit, add the private addend to the result modulo N and check status. The in-place addition is permitted here even though this public multiplier's output must be disjoint from its inputs.

Replace only the body of region \texttt{modproduct\_select}, retaining its BEGIN and END comments.

\Needspace{22\baselineskip}
\MARSCodeLabel{MIPS assembly}
\begin{minted}[samepage=true]{mars-idea}
mulmod_loop:
    ld $t0, 16($sp)
    ld $t1, 24($sp)
    or $t2, $t0, $t1
    beqz $t2, mulmod_done
    nop
    andi $t2, $t0, 1
    beqz $t2, mulmod_shift
    nop
    daddiu $a0, $sp, 32
    move $a1, $sp
    move $a2, $s0
    daddiu $a3, $sp, 32
    jal rsa_addmod128
    nop
    bnez $v0, mulmod_development
    nop
\end{minted}

Use recipe \texttt{s31-1} with the cumulative source through this edit. Stop at \texttt{mulmod\_shift} before continuing to the temporary or normal return.

\subsection{Expected observation}
The first selected addition produces 13 while retaining multiplier 11.

\Needspace{9\baselineskip}
\subsection{Check your understanding}
Why is the nested result update allowed in place?

\begin{itemize}
\item \textbf{A} The stack pointer makes all overlap harmless.
\item \textbf{B} The modular-addition contract explicitly permits that exact alias.
\item \textbf{C} Every public arithmetic output may overlap any input.
\end{itemize}
